% ============================================================================
% ISCHLSTROM Batteriemanagement (IBM) - Whitepaper
% Kompilieren: pdflatex ibm-whitepaper.tex  (zweimal fuer Verweise)
% ============================================================================
\documentclass[11pt,a4paper]{article}

\usepackage[T1]{fontenc}
\usepackage[utf8]{inputenc}
\usepackage[ngerman]{babel}
\usepackage{lmodern}
\usepackage{microtype}
\usepackage[a4paper,margin=2.5cm]{geometry}
\usepackage{parskip}
\usepackage{booktabs}
\usepackage{amsmath}
\usepackage{graphicx}
\usepackage{tikz}
\usetikzlibrary{positioning,arrows.meta,fit,backgrounds}
\usepackage[colorlinks=true,linkcolor=black,urlcolor=blue!50!black]{hyperref}

\newcommand{\code}[1]{\texttt{#1}}

\title{\textbf{ISCHLSTROM Batteriemanagement (IBM)}\\[0.5em]
\large Gemeinschaftsdienliche Steuerung privater Heimspeicher\\
in einer Erneuerbare-Energie-Gemeinschaft}
\author{Energiegemeinschaft ISCHLSTROM \\ \href{https://ischlstrom.org}{ischlstrom.org}}
\date{August 2026}

\begin{document}

\maketitle

\begin{abstract}
\noindent
Das ISCHLSTROM Batteriemanagement (IBM) steuert private Heimspeicher von
Mitgliedern einer Erneuerbare-Energie-Gemeinschaft (EEG) so, dass sie der
Gemeinschaft nützen, ohne dem einzelnen Haushalt zu schaden. Zwei Eingriffe
genügen dafür: An sonnigen Tagen regelt ein geschlossener Regelkreis die
Ladeleistung der Batterie so, dass sie den ganzen Tag gerade schnell genug
lädt, um am Abend voll zu sein; der restliche PV-Überschuss fließt laufend
in die Gemeinschaft, solange diese ihn braucht. In der Nacht speist die
Batterie kontrolliert in das Netz ein, solange die Gemeinschaft mehr
verbraucht als erzeugt. Beide
Eingriffe sind an Prognosen gebunden und mehrfach abgesichert: Ohne
verlässliche Daten passiert nichts, und der Wechselrichter fällt nach
wenigen Minuten von selbst in sein Werksverhalten zurück. Die Steuerung
lernt Batteriekapazität, Ladeleistung und nächtliche Hauslast jeder Anlage
selbstständig aus dem Betrieb; eine manuelle Parametrierung pro Anlage
entfällt. Dieses Papier beschreibt Architektur, Algorithmen,
Sicherheitsmechanismen und Betrieb des Systems im Detail.
\end{abstract}

\tableofcontents
\newpage

% ============================================================================
\section{Motivation}
% ============================================================================

Eine Erneuerbare-Energie-Gemeinschaft verteilt lokal erzeugten PV-Strom
unter ihren Mitgliedern. Ihr wirtschaftlicher und ökologischer Nutzen hängt
davon ab, wie gut Erzeugung und Verbrauch zeitlich zusammenpassen: Nur
Energie, die im selben Viertelstunden-Raster erzeugt und verbraucht wird,
zählt als gemeinschaftliche Deckung. ISCHLSTROM kauft als Gemeinschaft
niemals Strom zu; was die Mitglieder nicht aus der Gemeinschaft decken,
beziehen sie individuell von ihrem eigenen Lieferanten.

Private Heimspeicher verschärfen das Gleichzeitigkeitsproblem, statt es zu
lösen: Ab Sonnenaufgang lädt jede Batterie sofort, der Vormittags-Überschuss
der PV-Anlagen verschwindet in den Speichern, während die Gemeinschaft noch
Strom braucht. Am Abend und in der Nacht versorgt die volle Batterie nur den
eigenen Haushalt, während die Gemeinschaft aus dem Netz bezieht.

IBM dreht dieses Verhalten um, ohne den Eigennutzen des Speicherbesitzers zu
opfern:

\begin{itemize}
  \item \textbf{Tagsüber (Laderegelung):} Die Batterie lädt so langsam wie
    möglich, wird aber bis zum Abend trotzdem voll. Der nicht gespeicherte
    Überschuss deckt die Verbrauchsspitzen der Gemeinschaft direkt, vom
    frühen Vormittag an und über den ganzen Tag verteilt.
  \item \textbf{Nacht (forcierte Entladung):} Die Batterie speist einen
    berechneten Teil ihres Inhalts in das Netz ein, solange die Gemeinschaft
    mehr verbraucht als erzeugt. Eingespeist wird nur, was der kommende Tag
    laut Prognose sicher wieder in die Batterie lädt.
\end{itemize}

Geladen wird dabei ausschließlich aus der eigenen PV-Anlage, nie aus dem
Netz oder von anderen Mitgliedern.

% ============================================================================
\section{Systemarchitektur}
% ============================================================================

IBM besteht aus zwei Teilen: einem zentralen Server (die Website der
Gemeinschaft, \code{ischlstrom.org}) und je Anlage einem lokalen Gateway,
einem Raspberry~Pi mit openHAB (openHABian), der den Wechselrichter im
Heimnetz des Mitglieds ansteuert.

\begin{figure}[h]
\centering
\begin{tikzpicture}[
  node distance=0.9cm and 2.1cm,
  box/.style={draw, rounded corners, align=center, minimum height=2.4em,
              inner sep=6pt, font=\small},
  cloud/.style={box, fill=blue!8},
  pi/.style={box, fill=green!8},
  hw/.style={box, fill=orange!10},
  lbl/.style={font=\scriptsize, midway, align=center},
  arr/.style={-{Stealth}, thick}
]
  \node[cloud] (server) {Server \code{ischlstrom.org}\\Prognose, Ladefenster,\\Crossover, Wolken, Dashboard};
  \node[pi, below left=2.2cm and -1.2cm of server] (rules) {openHAB-Regeln\\(API-Abrufe, Status-Push)};
  \node[pi, right=of rules] (core) {Steuerungskern\\\code{core.js} (alle 5\,min)};
  \node[pi, right=of core] (adapter) {Adapter\\(je Hersteller)};
  \node[hw, right=1.7cm of adapter] (inverter) {Wechselrichter\\+ Batterie};

  \begin{scope}[on background layer]
    \node[draw, dashed, rounded corners, fit=(rules)(core)(adapter),
          label={[font=\scriptsize]below:{Raspberry Pi (openHABian) im Heimnetz des Mitglieds}}] {};
  \end{scope}

  \draw[arr] (server.south -| rules.north) -- (rules.north)
    node[lbl, left=1pt] {HTTPS\\(Abruf + Push)};
  \draw[arr] (rules) -- (core) node[lbl, above] {Items};
  \draw[arr] (core) -- (adapter) node[lbl, above] {3 Funktionen};
  \draw[arr] (adapter) -- (inverter) node[lbl, above] {Modbus /\\Binding};
\end{tikzpicture}
\caption{Gesamtarchitektur: Der Server liefert Prognosedaten, das lokale
Gateway entscheidet und steuert. Alle Verbindungen gehen vom Gateway aus;
am Router des Mitglieds wird nichts geöffnet.}
\label{fig:arch}
\end{figure}

\subsection{Rollenteilung Server und Gateway}

Der Server kennt die Gemeinschaft: die Viertelstunden-Prognose von Erzeugung
und Verbrauch, die Wetterdaten und die gemeldeten Kennwerte aller Anlagen.
Daraus berechnet er die Steuersignale (Abschnitt~\ref{sec:signale}). Das
Gateway kennt die Anlage: den Ladestand, den Wechselrichter und die lokal
gelernten Kennwerte. Die eigentliche Entscheidung, ob in diesem Moment
gesperrt oder entladen wird, fällt immer lokal auf dem Gateway. Fällt der
Server aus, passiert nichts Gefährliches: Ohne aktuelle Daten unterbleiben
die Eingriffe (Abschnitt~\ref{sec:sicherheit}).

\subsection{Kern und Adapter}

Die gesamte Entscheidungslogik liegt herstellerneutral in einem einzigen
Skript (\code{control/core.js}), das alle fünf Minuten läuft und pro Kunde
nie verändert wird. Alles Anlagenspezifische kommt aus openHAB-Items mit
Rückfallwerten im Code; fehlt ein Item, läuft die Steuerung mit der Vorgabe
weiter.

Herstellerspezifisch ist nur ein dünner \emph{Adapter}, den das Setup dem
Kern im selben Regel-Body voranstellt. Der Adapter-Kontrakt umfasst drei
Pflichtfunktionen und eine optionale:

\begin{center}
\small
\begin{tabular}{@{}lll@{}}
\toprule
Funktion & Semantik & Rückgabe \\
\midrule
\code{ibmReset()} & Wechselrichter sofort auf Werksverhalten & \code{\{ok\}} \\
\code{ibmPreventCharge(min)} & Batterieladen für \code{min} Minuten sperren & \code{\{ok\}} \\
\code{ibmForceDischarge(W, min)} & Entladung mit $\approx$\code{W} Watt erzwingen & \code{\{ok, appliedW?\}} \\
\code{ibmLimitCharge(W, min)} & optional: Ladeleistung auf $\approx$\code{W} Watt begrenzen & \code{\{ok, appliedW?\}} \\
\bottomrule
\end{tabular}
\end{center}

Jede Aktion gilt nur für einen Zeitschlitz von fünf Minuten und muss
hersteller\-seitig von selbst ablaufen (Schedule oder Revert-Timeout). Der
nächste Lauf des Kerns verlängert den Eingriff oder beendet ihn, indem er
schlicht nichts mehr kommandiert. \code{appliedW} meldet die nach
herstellerseitiger Quantisierung tatsächlich kommandierte Leistung (etwa
durch Prozent-Rundung); der Kapazitätsschätzer rechnet damit.
\code{ibmLimitCharge} definieren nur Profile, deren Hardware ein echtes
Ladelimit kennt, ohne dabei aus dem Netz laden zu können (SunSpec
\code{InWRte}, Victron DVCC); fehlt die Funktion, bildet der Kern die
Begrenzung per Puls-Weiten-Modulation über \code{ibmPreventCharge} nach
(Abschnitt~\ref{sec:ladesperre}).

Unterstützt werden derzeit fünf Profile: Fronius GEN24 (Binding-Actions),
Fronius Symo Hybrid (Modbus, SunSpec Model~124), Sigenergy SigenStor
(proprietäre Modbus-Register), Deye SG04/SG05LP3 (Modbus RTU hinter einem
RS485-Gateway, TOU-Fahrplan) und Victron Energy (Modbus TCP am GX-Gerät,
ESS-Register). Ein neuer Hersteller braucht nur ein neues Profilverzeichnis
mit Adapter; an den Setup-Skripten und am Kern ändert sich nichts.

\subsection{Zeitgesteuerte Regeln}

Auf dem Gateway laufen wenige, klar getrennte Regeln:

\begin{center}
\small
\begin{tabular}{@{}lll@{}}
\toprule
Regel & Aufgabe & Zeitplan \\
\midrule
\code{ibm\_cloud\_forecast} & Wolkenvorschau abrufen & stündlich \\
\code{ibm\_crossover} & Crossover-Zeiten abrufen & täglich 04:05 \\
\code{ibm\_ladesperre} & Ladefenster und Nachtbudget abrufen & stündlich \\
\code{ibm\_battery\_control} & Adapter + Kern: entscheiden und steuern & alle 5\,min \\
\code{ibm\_status\_push} & Zustand an das Dashboard melden & jede Minute \\
\code{ibm\_init} & Startwerte setzen, solange Items \code{NULL} sind & alle 10\,min \\
\code{ibm\_pause} & Pausentage herunterzählen & täglich 00:30 \\
\code{ibm\_watchdog} & Wechselrichter nach IP-Wechsel wiederfinden & bei Bedarf \\
\bottomrule
\end{tabular}
\end{center}

% ============================================================================
\section{Steuersignale des Servers}\label{sec:signale}
% ============================================================================

Der Server stellt vier Signale bereit. Die Gateways rufen sie über eine
schlanke HTTP-API ab; die Community-Signale sind öffentlich, die
individualisierten erfordern das Status-Token der Anlage.

\begin{center}
\small
\begin{tabular}{@{}llll@{}}
\toprule
Endpunkt & Methode & Inhalt & Zugriff \\
\midrule
\code{/api/eeginfo/ladefenster/v1} & GET & Community-Ladefenster & öffentlich \\
\code{/api/ibm/ladefenster/v1} & POST & Fenster + Nachtbudget + Ladefaktoren & Token \\
\code{/api/eeginfo/crossover/v1} & GET & Crossover-Zeiten der Woche & öffentlich \\
\code{/api/wolken/vorschau/v1} & GET & Bewölkung: Mittagsfenster + Stundenwerte & öffentlich \\
\code{/api/ibm/status/v1} & POST & Status-Push der Anlage & Token \\
\bottomrule
\end{tabular}
\end{center}

\subsection{Datengrundlage: die Tagesprognose}

Grundlage der Ladefenster-Signale ist die Viertelstunden-Prognose der
Gemeinschaft: ein statistisches Modell (Python) prognostiziert Erzeugung,
Verbrauch und Eigendeckung je Messpunkt aus Kalender-, Sonnenstands- und
Wetter-Merkmalen (Open-Meteo). Prognoseläufe werden versioniert gespeichert
und nie überschrieben; die API liest stets den neuesten Live-Lauf. Die
Wetterdaten selbst holt ein stündlicher Serverjob.

\subsection{Community-Ladefenster}

Aus dem Prognosetag berechnet der Server das Sperrfenster der Gemeinschaft.
Mit der prognostizierten Erzeugung $G_t$ und dem Verbrauch $V_t$ je
Viertelstunden-Slot, dem Überschuss $S_t = G_t - V_t$ und dessen
Tagesmaximum $S_{\max}$ gilt:

\begin{itemize}
  \item \textbf{Beginn} ist der erste Slot mit nennenswertem Sonnenschein,
    $G_t > 0{,}05 \cdot G_{\max}$.
  \item \textbf{Ende} ist der spätere Zeitpunkt von Vormittags-Crossover
    (erster Slot mit $G_t \geq V_t$) und dem ersten Slot mit
    $S_t \geq 0{,}75 \cdot S_{\max}$, gekappt am Slot des
    Überschussmaximums und um 14:00:
    \begin{equation}
      t_{\text{Ende}} \;=\; \min\Bigl(
        \max\bigl(t_{\text{Crossover}},\; t_{S \geq 0{,}75\,S_{\max}}\bigr),
        \; t_{S_{\max}},\; 14\!:\!00 \Bigr).
    \end{equation}
\end{itemize}

Die Batterien bleiben so bis in die Überschussspitze gesperrt, statt direkt
nach dem Crossover loszuladen. Erwartet die Prognose keinen Überschuss,
liefert die API kein Ende und es wird nicht gesperrt. Jede Antwort trägt
das Datum des Prognosetags; die Gateways verwerfen Fenster mit fremdem
Datum.

\subsection{Individualisiertes Sperr-Ende}\label{sec:individuell}

Für Anlagen mit Status-Token rechnet der Server das Fensterende je Anlage,
aus deren zuletzt gemeldeten Kennwerten (Kapazität $C$, Ladeleistung
$P_{\text{Lade}}$, Abschnitt~\ref{sec:lernen}). Die Idee: Die erreichbare
Ladeleistung folgt dem Erzeugungsprofil des Tages. Der Server normiert das
Profil auf den Maximalslot des gesamten Prognoselaufs (nicht des Tages,
sonst sähe ein trüber Tag wie ein klarer aus) und integriert rückwärts von
der Abend-Deadline $t_D$ (Ende des letzten Überschuss-Slots minus
60~Minuten):
\begin{equation}
  E(t) \;=\; \sum_{s=t}^{t_D} P_{\text{Lade}} \cdot
  \frac{G_s}{G_{\max}^{\text{Lauf}}} \cdot 0{,}25\,\text{h}.
\end{equation}
Das Sperr-Ende ist der späteste Slot $t$, ab dem
$E(t) \geq 0{,}95 \cdot C \cdot 1{,}3$ noch geladen wird (95\% der
Kapazität mit Sicherheitsfaktor 1,3), gekappt um 14:00. Reicht der ganze
Tag nicht, wird nicht gesperrt. Individualisiert wird nur bei plausiblen
Kennwerten (Kapazität 1 bis 100\,kWh, Ladeleistung 0,3 bis 30\,kW), und
das Ende muss nach dem Fensterbeginn und nicht vor 05:00 liegen. Die
Antwort ist als individuell markiert; das Gateway übernimmt sie dann
unverändert und lässt die eigene Rechnung nach
Gleichung~\eqref{eq:localend} aus. Greift auf dem Gateway die
Laderegelung (Abschnitt~\ref{sec:ladesperre}), spielt das Ende keine Rolle
mehr; es bleibt der Rückfall für Anlagen ohne belastbare Schätzwerte.

\subsection{Nachtbudget}\label{sec:nachtbudget}

Mit dem individualisierten Ladefenster liefert der Server das
Entladebudget der kommenden Nacht. Er summiert zunächst, was die Anlage in
den nächsten 24~Stunden laut Prognoseprofil laden kann (dieselbe
Integration wie in Abschnitt~\ref{sec:individuell}, über alle Slots der
nächsten 24~Stunden), zieht eine Eigenbedarfsreserve von 24~Stunden
gelernter Hauslast $P_{\text{Haus}}$ ab und behält einen Abschlag von 20\%
ein:
\begin{equation}
  B \;=\; \max\Bigl(0,\; E_{\text{ladbar,24h}}
    - \frac{P_{\text{Haus}}}{1000} \cdot 24\,\text{h}\Bigr) \cdot 0{,}8.
\end{equation}
Meldet die Anlage keine plausible Hauslast (50 bis 3000\,W), rechnet der
Server mit 500\,W. Bei mehrtägigem Schlechtwetter ist das Budget schlicht
null: Die Batterie behält ihren Inhalt für das eigene Haus. Wie das
Gateway das Budget in einen Ziel-Ladestand übersetzt, zeigt
Gleichung~\eqref{eq:ziel}.

\subsection{Crossover-Zeiten}

Die Crossover-Zeiten sind die typischen Uhrzeiten, zu denen die
gemeinschaftliche Erzeugung den Gesamtverbrauch morgens über- und abends
unterschreitet. Sie stammen aus den historischen Viertelstunden-Messdaten
der Gemeinschaft: je Tag des Jahres über alle Jahre gemittelt und zu
Kalenderwochen zusammengefasst (materialisierte Sicht, monatlich
aufgefrischt). Anders als das Ladefenster sind sie also kein
Prognoseprodukt, sondern ein saisonales Klimamittel; sie ändern sich
langsam und sind deshalb als Wochenwert ausreichend genau.

\subsection{Wolkenvorschau und Stundenwerte}

Die Wolkenvorschau ist der mittlere Bewölkungsgrad über dem nächsten
Mittagsfenster von 10:00 bis 14:00 (nach 12:00 zählt bereits das Fenster
von morgen), berechnet aus den stündlich importierten
Open-Meteo-Vorhersagedaten. Sie steuert die nächtliche Entladung, die
genau diese Frage stellt: Wird die Batterie morgen wieder voll? Fehlen
Wetterdaten, antwortet die API bewusst mit einem Fehler statt mit 0, denn
0\% würde von den Gateways als voller Sonnenschein gelesen.

Für die Laderegelung, die bis in den Abend läuft, ist ein einzelner
Mittagswert zu grob und ab 12:00 sogar der falsche Tag. Dieselbe API
liefert deshalb zusätzlich die \emph{stündlichen} Bewölkungswerte vom
Beginn der laufenden Stunde bis Mitternacht, samt Gültigkeitsdatum. Das
Gateway legt sie als JSON ab; die Regelung gewichtet damit jede
verbleibende Stunde einzeln (Abschnitt~\ref{sec:ladesperre}).

\subsection{Ladefaktoren}\label{sec:ladefaktoren}

Mit dem individualisierten Ladefenster liefert die Token-API außerdem die
stündlichen \emph{Ladefaktoren} des heutigen Tages: je Stunde das Mittel
der normierten Erzeugung $G_s / G_{\max}^{\text{Lauf}}$ über die vier
Viertelstunden-Slots (dieselbe Normierung wie in
Abschnitt~\ref{sec:individuell}), dazu die Abend-Deadline. Ein Faktor von
0,6 heißt: In dieser Stunde lädt die Anlage voraussichtlich mit 60\% ihrer
gelernten Spitzenrate. Anders als die Bewölkung enthalten die Faktoren
auch den Sonnenstand, sie sind also das physikalisch genauere Gewicht für
die Restladezeit der Laderegelung.

% ============================================================================
\section{Teil A: Laderegelung am Tag}\label{sec:ladesperre}
% ============================================================================

\subsection{Grundprinzip: Regelkreis statt Sperrfenster}

An sonnigen Tagen wird die Ladeleistung der Batterie dynamisch geregelt.
In jedem Fünf-Minuten-Zyklus berechnet der Kern die Ziel-Ladeleistung neu:
fehlende Energie geteilt durch die \emph{effektive Restladezeit} bis zur
Abend-Deadline $t_D$ (Abend-Crossover minus 60~Minuten),
\begin{equation}
  P_{\text{Soll}}(t) \;=\;
  \frac{(0{,}95 - \text{SoC}(t)) \cdot C}{T_{\text{eff}}(t)} \cdot 1{,}1,
  \qquad
  T_{\text{eff}}(t) \;=\; \sum_{h\,\in\,[t,\,t_D)} f_h \cdot \Delta t_h
  \label{eq:psoll}
\end{equation}
mit Kapazität $C$ in kWh und Live-Ladestand $\text{SoC}(t)$ als Anteil.
Die Restzeit ist sonnengewichtet: Jede verbleibende Stunde $h$ zählt nur
mit ihrem erwarteten Ertragsfaktor $f_h$ und ihrer Überlappung
$\Delta t_h$ mit $[t, t_D)$. Der Faktor kommt in absteigender Priorität
aus den stündlichen Ladefaktoren der Token-API
($f_h = G_h / G_{\max}^{\text{Lauf}}$, Abschnitt~\ref{sec:ladefaktoren};
physikalisch exakt inklusive Sonnenstand), sonst aus den stündlichen
Bewölkungswerten der Wolken-API ($f_h = 1 - w_h/100$). Nur wenn beides
fehlt (älterer Server), zählt jede Stunde gleich
($T_{\text{eff}} = t_D - t$) und statt des festen Aufschlags 1,1 gleicht
ein wolkenabhängiger Sicherheitsfaktor $1{,}1 + w/100 \cdot 0{,}5$ den
Nachmittagsabfall pauschal aus.

Die Batterie lädt so den ganzen Tag gerade schnell genug, um am Abend voll
zu sein; der gesamte restliche PV-Überschuss fließt laufend in die
Gemeinschaft. Weil auf den Live-Ladestand geregelt wird, ist der Kreis
geschlossen: Zieht es zu, bleibt der Ladestand zurück, $P_{\text{Soll}}$
steigt und die Begrenzung löst sich von selbst; wird es sonniger als
vorhergesagt, sinkt $P_{\text{Soll}}$ und mehr geht ins Netz.
Prognosefehler regeln sich so alle fünf Minuten aus, statt einmal am
Vormittag verwettet zu werden. Die aktuelle effektive Restladezeit zeigt
das Gateway auf seiner Bedienseite an und meldet sie per Status-Push an
das Vorstands-Dashboard.

Geregelt wird nur zwischen dem Fensterbeginn des Servers (erster
Sonnenschein der Tagesprognose, gültig nur für das mitgelieferte Datum) und
der Abend-Deadline, und nur bei frischer, sonniger Vorschau: Liegt die
Bewölkung über der Schwelle (Vorgabe 75\%), bleibt das Laden frei, denn an
einem trüben Tag würde die Batterie sonst womöglich gar nicht mehr voll.
Der Wächter verwendet dafür die mittlere Bewölkung der Reststunden bis zur
Deadline; nur ohne Stundendaten die Mittagsfenster-Vorschau (die ab 12:00
bereits den morgigen Tag meint).

\subsection{Umsetzung: direktes Limit oder Puls-Weiten-Modulation}

Definiert der Adapter \code{ibmLimitCharge}, wird $P_{\text{Soll}}$ direkt
kommandiert (quantisiert auf 100-W-Schritte). Sonst bildet der Kern die
Begrenzung per Puls-Weiten-Modulation über \code{ibmPreventCharge} nach:
Der Sperranteil
\begin{equation}
  d \;=\; 1 - \frac{P_{\text{Soll}}}{P_{\text{Lade}}}
  \label{eq:duty}
\end{equation}
(mit der gelernten Ladeleistung $P_{\text{Lade}}$,
Abschnitt~\ref{sec:lernen}) bestimmt, welcher Anteil der Zeit gesperrt
wird. Entschieden wird je 15-Minuten-Block (drei Slots): Ein
Bresenham-Akkumulator sammelt je Slot den Sperranteil als
\emph{Sperrschuld} an, jeder gesperrte Slot zieht 1 ab; im Mittel entsteht
so genau die Ziel-Leistung. Eine Hysterese auf der Sperrschuld (Sperren ab
2,0, Freigeben unter 1,0) verhindert Flattern an den Blockgrenzen, und die
Batterie schaltet nie schneller als im 15-Minuten-Raster. Zwei Deckel
sichern die Regelung ab: Unter einem Sperranteil von 10\% wird gar nicht
begrenzt, und über 90\% wird gekappt; ein Schätzfehler kann das Laden also
nie ganz würgen. Fällt die Begrenzung ohnehin über die gelernte
Ladeleistung hinaus ($P_{\text{Soll}} \geq P_{\text{Lade}}$), lädt die
Batterie unbegrenzt.

\subsection{Rückfall: das Sperrfenster}\label{sec:sperrfenster}

Die Regelung braucht belastbare Schätzwerte für Kapazität und
Ladeleistung. Solange sie fehlen (frisch installierte Anlage) oder die
Regelung abgeschaltet ist, gilt das klassische Sperrfenster: Das Laden wird
in Fünf-Minuten-Schritten komplett gesperrt, vom ersten Sonnenschein bis zu
einem berechneten Sperr-Ende, danach lädt die Batterie mit voller Leistung.
Für das Sperr-Ende gibt es drei Stufen, in absteigender Priorität:

\begin{enumerate}
  \item \textbf{Individualisiertes Server-Ende:} Anlagen mit Status-Token
    rufen das Ladefenster über die Token-API ab. Der Server berechnet das
    Ende dann je Anlage aus dem Erzeugungsprofil des Prognosetags und den
    zuletzt gemeldeten Kennwerten (Kapazität, Ladeleistung): so spät, dass
    die Batterie am Abend gerade noch voll wird
    (Abschnitt~\ref{sec:individuell}). Die Antwort ist als individuell
    markiert; der Kern übernimmt sie unverändert.
  \item \textbf{Lokal berechnetes Ende:} Ohne individualisierte Antwort
    rechnet der Kern selbst, sobald Kapazität und Ladeleistung gelernt sind
    (Abschnitt~\ref{sec:lernen}). Rückwärts vom Abend:
    \begin{equation}
      t_{\text{Ende}} \;=\; \min\Bigl(
        t_{\text{Abend-Crossover}} - 60\,\text{min}
        \;-\; \underbrace{\frac{0{,}95 \cdot C}{P_{\text{Lade}}} \cdot 1{,}3 \cdot 60\,\text{min}}_{\text{Ladezeit mit Sicherheitsaufschlag}},
        \;\; 13\!:\!00\Bigr)
      \label{eq:localend}
    \end{equation}
    mit Kapazität $C$ in kWh und gelernter Ladeleistung $P_{\text{Lade}}$
    in kW. Als fehlende Energie gelten fest 95\% der Kapazität, nicht der
    Live-Ladestand: Das Ergebnis soll den ganzen Vormittag konstant sein
    und nicht pendeln, sobald das Laden den Ladestand hebt. Liegt das
    berechnete Ende vor dem Fensterbeginn, braucht die Anlage den ganzen
    Tag zum Laden und es wird gar nicht gesperrt.
  \item \textbf{Community-Ende des Servers:} Solange die Schätzer noch
    nicht belastbar sind, gilt das für die ganze Gemeinschaft berechnete
    Fensterende unverändert.
\end{enumerate}

Greift die Laderegelung, sind alle drei Stufen gegenstandslos: Der
geschlossene Regelkreis reagiert auf den Live-Ladestand besser, als es jede
Vorausberechnung eines Endes kann. Der Fensterbeginn (erster Sonnenschein)
und die Datumsprüfung kommen in jedem Fall vom Server.
Plausibilitätsfenster verwerfen Unsinniges: Der Beginn muss zwischen 4 und
12~Uhr liegen, das Ende zwischen 5 und 15~Uhr.

\subsection{Wirkung für die Gemeinschaft}

Die Verbrauchsspitzen der Gemeinschaft werden direkt aus der PV gedeckt,
und die Einspeisung verteilt sich über den ganzen Tag statt erst nach einem
Sperr-Ende einzusetzen. Das hält die Gemeinschaft nach dem Crossover im
Plus, fängt Abregelungsverluste und verkürzt nebenbei die Standzeit der
Batterie bei 100\% Ladung.

% ============================================================================
\section{Teil B: Forcierte Entladung in der Nacht}\label{sec:entladung}
% ============================================================================

\subsection{Entladefenster}

Entladen wird vom abendlichen bis zum morgendlichen Crossover, also genau
solange die Gemeinschaft mehr verbraucht als erzeugt. Die Crossover-Zeiten
holt das Gateway täglich vom Server. Liegen keine plausiblen Zeiten vor
(Server nie erreichbar gewesen oder Werte unbrauchbar), wird nicht entladen;
ein Ersatzfenster gibt es bewusst nicht.

\subsection{Entladeleistung aus der Wolkenvorschau}

Die Entladeleistung hängt von der Bewölkungsvorschau für das nächste
Sonnenfenster ab, linear interpoliert zwischen einem Minimum und einem
Maximum:
\begin{equation}
  P_{\text{Entladung}} \;=\; P_{\max} - \frac{w}{100}\,(P_{\max} - P_{\min}),
  \qquad w = \text{Bewölkung in \%}.
  \label{eq:power}
\end{equation}
Bei klarer Vorschau ($w=0$) wird mit $P_{\max}$ entladen, weil die Batterie
am nächsten Tag sicher wieder voll wird; bei trüber Vorschau entsprechend
weniger. Zwei Sonderfälle:

\begin{itemize}
  \item \textbf{Trüb-Stopp:} Liegt die Vorschau über der Wolkenschwelle
    (dieselbe wie bei der Ladesperre), wird gar nicht eingespeist. Über
    eine lange Nacht entlädt auch minimale Leistung die Batterie fast
    vollständig; am trüben Folgetag müsste das Mitglied dann selbst Strom
    zukaufen. Es gilt die Symmetrie: Nur wenn morgen früh gesperrt würde
    (sonnig), wird heute Nacht auch eingespeist.
  \item \textbf{Ausfall der Vorschau:} Ist die Wolkenvorschau älter als
    drei Stunden oder ungültig, wird konservativ so entladen, als wäre der
    nächste Tag komplett bewölkt, also mit $P_{\min}$. Die Batterie läuft
    bei einem API-Ausfall damit nie mit Maximalleistung leer.
\end{itemize}

\subsection{Dynamische Leistungsgrenzen (C-Rate)}

$P_{\min}$ und $P_{\max}$ sind einstellbar (Vorgabe 1000 und 3000\,W),
werden aber ersetzt, sobald die Kapazitätsschätzung belastbar ist
(Abschnitt~\ref{sec:lernen}): Dann gilt eine C-Raten-Spanne von 0,1\,C bis
0,3\,C. Eine 10-kWh-Batterie speist also mit 1000 bis 3000\,W ein, eine
5-kWh-Batterie mit 500 bis 1500\,W. Unabhängig von Einstellung und
Schätzung gilt eine harte Sicherheits-Obergrenze von 5000\,W, die nie
überschritten wird.

\subsection{Untergrenzen: Reserve und Nachtbudget}

Zwei Grenzen schützen den Haushalt vor einer leeren Batterie:

\begin{itemize}
  \item \textbf{Reserve:} Unter den eingestellten Mindest-Ladestand
    (\code{IBM\_MIN\_BATTERY\_CHARGE}, plausibel zwischen 5 und 90\%) wird
    nie forciert entladen.
  \item \textbf{Nachtbudget:} Die Token-API liefert mit dem Ladefenster ein
    Entladebudget $B$ in kWh: so viel darf heute Nacht eingespeist werden,
    ohne dass die Batterie am Folgetag dem eigenen Haus fehlt. Beim ersten
    Entlade-Zyklus der Nacht wird daraus ein Ziel-Ladestand berechnet und
    für die ganze Nacht festgehalten:
    \begin{equation}
      \text{SoC}_{\text{Ziel}} \;=\; \max\Bigl(\text{SoC}_{\min},\;
        \text{SoC}_{\text{Abend}} - \frac{B}{C}\cdot 100\Bigr).
      \label{eq:ziel}
    \end{equation}
    Das Festhalten verhindert, dass ein später aktualisiertes Budget in
    derselben Nacht doppelt ausgegeben wird; eine Lücke von mehr als
    60~Minuten gilt als neue Nacht. Ohne (aktuelles) Budget, etwa ohne
    Token oder bei veraltetem Abruf, gilt wie zuvor nur die Reserve.
\end{itemize}

Nach dem Entlade-Stopp versorgt der Wechselrichter das Haus weiter aus der
Batterie, bis zu seiner eigenen Reserve von wenigen Prozent; genau dieser
Abschnitt liefert die Hauslast-Messung (Abschnitt~\ref{sec:lernen}).

% ============================================================================
\section{Selbstlernende Anlagenkennwerte}\label{sec:lernen}
% ============================================================================

Bei der Einrichtung sind weder Batteriekapazität noch reale Ladeleistung
noch nächtliche Hauslast bekannt, und niemand soll sie pro Anlage pflegen
müssen. Die Steuerung schätzt alle drei Werte aus dem laufenden Betrieb.
Alle Schätzer teilen dasselbe Muster: Messstrecken über mehrere
Fünf-Minuten-Läufe, Plausibilitätsfenster je Stichprobe, Neuaufsetzen nach
Lücken oder widersprüchlichen Ladestandsbewegungen, gleitende Mittelung und
eine Mindestzahl von Stichproben, bevor der Wert verwendet wird. Der
Zustand jedes Schätzers liegt als JSON in einem persistierten Item und
überlebt Neustarts.

\subsection{Batteriekapazität}

Während der forcierten Entladung ist die Batterieleistung bekannt, denn sie
wird kommandiert. Der Kern summiert die entnommene Energie über die
Fünf-Minuten-Läufe auf; fällt der Ladestand um mindestens 8~Prozentpunkte,
ergibt sich eine Stichprobe
\begin{equation}
  C_{\text{Stichprobe}} \;=\;
  \frac{E_{\text{entnommen}}\,[\text{kWh}]}{\Delta\text{SoC}\,[\%]}\cdot 100,
\end{equation}
die mit Gewicht 0,3 gleitend in die Schätzung einfließt. Stichproben unter
1 oder über 100\,kWh werden verworfen; verwendet wird die Schätzung ab drei
akzeptierten Stichproben. Eine typische Nacht liefert mehrere Stichproben,
die Schätzung ist also meist nach der ersten Nacht belastbar. Zieht der
Haushalt nachts mehr als kommandiert, fällt die Schätzung etwas zu klein
aus; das ist gewollt konservativ, denn die daraus abgeleitete
Entladeleistung ist dann eher zu niedrig als zu hoch.

\subsection{Ladeleistung}

Gemessen wird nur, wenn die Messung die real erreichbare Ladeleistung
widerspiegelt: tagsüber zwischen Fensterbeginn und Abend-Deadline, wenn die
Batterie frei laden darf, die Vorschau Sonne meldet und der Ladestand unter
95\% liegt (darüber drosselt der Wechselrichter selbst). In freien
PWM-Blöcken der Laderegelung wird weiter gemessen (dort lädt die Batterie
unbegrenzt); unter einem direkten Ladelimit nie, denn die Stichprobe würde
das Limit statt der Anlage messen und die Schätzung nach unten ziehen.
Steigt der Ladestand um mindestens 5~Prozentpunkte, ergibt geladene Energie
(aus der Kapazität) geteilt durch die Messdauer eine Stichprobe in kW.

Die Mittelung ist bewusst \emph{asymmetrisch und dauergewichtet}: Eine
Stichprobe unter der bisherigen Schätzung (etwa Dunst trotz sonniger
Vorschau) zählt mit Gewicht 0,5 je Messstunde, eine darüber nur mit 0,15 je
Messstunde, gedeckelt bei 0,6. Die Schätzung liegt damit nahe an der
unteren Hüllkurve der letzten Tage. Das ist die sichere Seite: Die Batterie
muss abends voll sein, also wird die Ladezeit eher über- als unterschätzt
und die Sperre endet eher zu früh als zu spät. Die Dauergewichtung
verhindert zugleich, dass schnelle Mittagsphasen mit vielen Stichproben je
Stunde den Schnitt systematisch nach oben ziehen.

\subsection{Nächtliche Hauslast}

Nach dem Entlade-Stopp versorgt die Batterie allein das Haus. Der
Ladestandsabfall klar unterhalb der Entlade-Reserve (und oberhalb der
wechselrichtereigenen Reserve) ergibt daher direkt die Hauslast:
Stichproben zwischen 50 und 3000\,W fließen mit Gewicht 0,3 in die
Schätzung ein. Der Wert geht per Status-Push an den Server, der daraus die
Eigenbedarfsreserve des Nachtbudgets ableitet; der Kreis zwischen lokalem
Lernen und serverseitiger Individualisierung schließt sich hier.

% ============================================================================
\section{Sicherheit und Robustheit}\label{sec:sicherheit}
% ============================================================================

Das System greift in fremde Hardware ein und läuft unbeaufsichtigt in
Privathaushalten. Entsprechend defensiv ist es ausgelegt; die Grundregel
lautet: \emph{Im Zweifel nichts tun.} Der Werkszustand des Wechselrichters,
also Eigenverbrauchsoptimierung, ist immer der sichere Rückfall.

\begin{itemize}
  \item \textbf{Fail-Safe im Wechselrichter:} Jeder Eingriff gilt für fünf
    Minuten und läuft herstellerseitig von selbst ab. Stirbt der Pi, hängt
    openHAB oder fällt das Netzwerk aus, arbeitet der Wechselrichter nach
    spätestens fünf Minuten wieder ab Werk. Kann ein Hersteller das nicht
    garantieren, dokumentiert sein Profil das Restrisiko ausdrücklich.
  \item \textbf{Kein Eingriff ohne Daten:} Ohne plausible Crossover-Zeiten
    keine Entladung; ohne gültiges, datumsrichtiges Ladefenster weder
    Laderegelung noch Sperre; ohne frische Wolkenvorschau (älter als drei
    Stunden) keine Ladebegrenzung und nur minimale Entladung.
  \item \textbf{Laden nie ganz würgen:} Die Laderegelung begrenzt höchstens
    auf 90\% Sperranteil bzw. mindestens 10\% der gelernten Ladeleistung.
    Selbst ein grober Schätzfehler lässt die Batterie weiter laden, und der
    Regelkreis korrigiert ihn über den Live-Ladestand.
  \item \textbf{Plausibilitätsfenster überall:} Jede Konfigurations- und
    Messgröße wird gegen einen erlaubten Bereich geprüft; Unplausibles wird
    verworfen und durch den Rückfallwert ersetzt, mit Logmeldung.
  \item \textbf{Harte Leistungsgrenze:} 5000\,W Entladeleistung werden nie
    überschritten, weder durch Einstellungen noch durch die
    Kapazitätsschätzung.
  \item \textbf{Untergrenzen des Ladestands:} Reserve und Nachtbudget
    (Abschnitt~\ref{sec:entladung}) verhindern, dass der Haushalt mit
    leerer Batterie in einen trüben Tag startet.
  \item \textbf{Hauptschalter und Pause:} Der Hauptschalter steht nach der
    Installation auf \code{OFF}; erst das Mitglied schaltet die Steuerung
    scharf. Bei \code{ON} wird in jedem Zyklus zuerst
    \code{ibmReset()} gerufen, der Sollzustand ist also selbstheilend.
    Über die Pausenfunktion setzt das Mitglied die Steuerung für eine
    gewählte Zahl von Tagen aus (etwa bei Urlaub oder Handwerkerterminen);
    ein nächtlicher Zähler zählt herunter, danach läuft alles von selbst
    wieder an.
  \item \textbf{Robust gegen fehlende Konfiguration:} Jedes fehlende oder
    \code{NULL}-Item wird durch einen dokumentierten Rückfallwert ersetzt;
    die Steuerung läuft auch unvollständig eingerichtet weiter.
\end{itemize}

% ============================================================================
\section{Betrieb und Überwachung}
% ============================================================================

\subsection{Status-Push und Vorstands-Dashboard}

Auf Wunsch meldet jede Anlage minütlich ihren Zustand an den Server:
Ladestand, Wechselrichter-Status, Batterie-, Netz-, PV- und Lastleistung,
alle Schalterstellungen, die gelernten Kennwerte, die zuletzt empfangenen
Steuersignale sowie Ziel-Ladeleistung und effektive Restladezeit der
Laderegelung. Alle fünf Minuten geht zusätzlich der volle Zustand mit: die
Warnungen und Fehler aus dem openHAB-Log der letzten 24~Stunden,
Versionsstände (IBM-Paket, openHAB, Java, Betriebssystem), ausstehende
Betriebssystem-Updates und der Systemzustand des Pi (CPU-Temperatur,
Unterspannungs-Register, SD-Karten-Füllstand, RAM/Swap, Boot-Zeitpunkt);
diese Sammler starten eigene Prozesse und laufen deshalb nicht minütlich.
Der Server mischt schlanke Meldungen in den letzten Stand und legt nur
volle Meldungen im Verlauf ab. Übertragen werden ausschließlich diese
Betriebsdaten, keine Verbrauchsdaten des Haushalts.

Der Vorstand sieht alle Anlagen live auf einem Dashboard
(\code{/board/openhab}), inklusive Wochendiagrammen von Ladestand,
Batterieleistung und Netzeinspeisung aus der Batterie. Ausbleibende
Meldungen färben die Anlage nach 15~Minuten gelb, nach einer Stunde rot;
das Dashboard ist damit zugleich die Ausfallüberwachung der Flotte.

Die Netzeinspeisung aus der Batterie ist dabei eine berechnete Größe: Von
der Entladeleistung fließt nur der Teil ins Netz, den das Netz gerade
aufnimmt, der Rest versorgt den Haushalt:
\begin{equation}
  P_{\text{Batterie}\rightarrow\text{Netz}} \;=\;
  \min\bigl(\max(P_{\text{Batterie}},0),\; \max(-P_{\text{Netz}},0)\bigr)
\end{equation}
mit der Vorzeichenkonvention Batterie positiv beim Entladen, Netz negativ
bei Einspeisung.

\subsection{Authentifizierung}

Jede Anlage authentifiziert sich mit einem Status-Token, das der Vorstand
vor der Installation auf dem Dashboard je Mitglied erzeugt und das bei der
Einrichtung in der Konfiguration des Gateways landet. Dasselbe Token
individualisiert den Ladefenster-Abruf; beide Token-Endpunkte sind bewusst
POST-Aufrufe mit dem Token im Body, damit es in keinem Access-Log
auftaucht. Meldungen mit unbekanntem Token weist der Server ab; das
Löschen des Tokens auf dem Dashboard widerruft den Zugang sofort. Auf dem
Server ist der letzte gemeldete Zustand je Anlage gespeichert (begrenzt
auf 16\,kB), dazu eine Verlaufshistorie im Fünf-Minuten-Raster mit
30~Tagen Aufbewahrung, aus der die Dashboard-Diagramme entstehen; die
Logauszüge gehen nicht in die Historie ein. Es werden keine Client-IPs
und keine personenbezogenen Daten protokolliert.

\subsection{Watchdog und Fernwartung}

Vergibt der Router dem Wechselrichter per DHCP eine neue IP, findet ein
Watchdog das Gerät im lokalen /24-Netz wieder, verifiziert es über die
Seriennummer (nie wird ein fremdes Gerät übernommen) und trägt die neue
Adresse per REST-API in das Thing ein. Für die Wartung hält jeder Pi einen
ausgehenden WireGuard-Tunnel zum Wartungsserver offen; am Router des
Mitglieds muss nichts geöffnet werden, und durch den Tunnel läuft
ausschließlich das Wartungsnetz. Sicherheitsupdates des Betriebssystems
spielt die Anlage über \code{unattended-upgrades} automatisch ein.

\subsection{Installation}

Die Einrichtung beim Mitglied besteht aus dem Flashen des
openHABian-Images, dem Anlegen des Admin-Kontos und einem einzigen
Kommando:
\begin{center}
\code{curl -fsSL https://ischlstrom.org/ibm/install.sh | sudo bash}
\end{center}
Der Assistent findet den Wechselrichter im Netz, legt Things und Items an,
installiert Regeln und Bedienoberfläche und verifiziert das Ergebnis. Alle
Schritte sind idempotent; ein Paket-Update übernimmt die bestehende
Konfiguration. Bedient wird die Anlage danach über die openHAB Main UI am
Handy: Hauptschalter, Teilfunktionen, Pause und eine Expertenseite.

% ============================================================================
\section{Zusammenfassung}
% ============================================================================

IBM zeigt, dass sich private Heimspeicher mit geringem Aufwand
gemeinschaftsdienlich steuern lassen, ohne Komfort- oder Kostennachteil für
den Besitzer. Die Architektur trennt strikt: Der Server rechnet mit dem
Wissen der Gemeinschaft, das Gateway entscheidet mit dem Wissen der Anlage,
und der Wechselrichter fällt ohne frische Kommandos von selbst in sein
Werksverhalten zurück. Selbstlernende Schätzer für Kapazität, Ladeleistung
und Hauslast ersetzen jede manuelle Parametrierung; das Kern-Adapter-Muster
hält den Aufwand für neue Wechselrichter-Hersteller bei wenigen Dutzend
Zeilen. Nach demselben Muster lässt sich das System auf weitere
Gemeinschaften übertragen.

\end{document}
